\minitopic{专题研讨：语境、断言与行动的三角关系}银行案例之所以经久不衰，不是因为它直接证明了某一理论，而是因为同一组直觉至少容纳三种解释。第一种把变化放在“知道”这个词的适用标准上：低风险谈话中，普通反差可能性被忽略；高风险谈话中，星期六不开门等可能性被提升，于是归属句“她知道银行明天营业”的真值随归属者语境变化。第二种保持知识条件稳定，把变化理解为行动要求：她一直知道营业时间，但在房屋交易即将失败时，仍有理由复核。第三种让实践利害进入知识本身：当错误代价升高，相同证据不再足以使主体知道。语境主义、纯粹行动解释和实践侵入因此构成需要逐项比较的三角，而不是“标准会不会提高”这一道选择题。\parencite{derose1992,hawthorne2004,stanley2005}

\minitopic{先固定谁的语境}讨论中最常见的混乱，是没有区分认识主体、知识归属者和评价者。假设林在周五看见银行公告，周六不在场；甲在普通聊天中说“林知道银行周六营业”，乙在法庭上讨论林是否尽到查验义务。语境主义主要让甲、乙各自谈话中的标准发挥作用，即使林的证据和风险完全不变，两个归属句也可能得到不同评价。主体敏感论则把林本人的利害或注意状态写进知识条件。实践侵入论通常关心林在当时作出何种决定以及错误代价多大。一个案例若同时改变三者，就无法辨认究竟是哪一种机制产生直觉。 设计思想实验时应采用“单变量改写”。先保持主体证据、行动和风险不变，只改变归属者是否提到怀疑可能；再保持谈话不变，只提高主体决策代价；最后保持风险不变，让主体亲自注意到一个反例。每次改写都记录四个判断：“知道”是否为真、主体是否应当相信、是否可以断言、是否可以据此行动。若四项判断始终一起变化，证据支持统一理论；若它们分离，就需要多种规范。

\minitopic{断言规范能够做多少工作}日常争论常从一句话的适当性推回知识条件。若一位医生说“我知道病变是良性的”，听者通常理解为医生愿意承担一定的证据责任。知识规范论把知识视为断言的适当标准：断言 \(p\) 至少要求说话者知道 \(p\)。然而不适当断言未必表明缺乏知识。会议时间虽已知道，在对方即将购买昂贵机票时，补一句“我再核对一下”可能是礼貌、风险沟通或专业职责，而不是撤回知识。相反，一个人也可能在压力下不当地断言自己并不知道的事。 可以用“撤回—修正—保留”测试区分这些现象。面对风险提醒，说话者若改口“那么我其实不知道”，似乎接受了知识标准变化；若说“我知道，但还不足以据此完成手术”，则把知识与行动分开；若说“我原本知道，新的反证使我现在不知道”，变化来自证据击败，而非单纯利害。研究者不应只问受试者是否同意原句，还应观察他们选择哪一种解释性续句。 断言还受到信息量和角色义务限制。护士即使知道药物通常安全，也不应省略患者已知过敏这一相关限定；工程师知道桥梁符合常规荷载，也不能在讨论超限运输时只说“桥是安全的”。这些例子说明“可以断言什么”同时取决于知识、相关性、听者可能误解和制度角色。若把一切断言缺陷都归因于不知道，知识概念就会吸收本应由交往伦理解释的现象。

\minitopic{行动理由与高风险证据}实践侵入最强的动机来自知识与理性行动之间的联系。若林知道银行周六营业，而去银行是她避免违约的唯一方案，似乎她可以把营业作为行动理由。高风险条件下我们却要求更强查验，于是有人推论她尚不知道。但这一步依赖一项桥梁原则：“只要知道 \(p\)，就可以在任何相关行动中无条件以 \(p\) 为前提。”桥梁原则过强时，任何极小错误概率都会在巨大损失下取消知识；过弱时，实践利害又无法进入知识。 更细的模型把行动前提分层。知识可以使 \(p\) 成为一个合格理由，却不保证它压倒所有其他理由；复核成本、损失规模、可逆性和时间压力共同决定行动。飞行员可以知道仪表读数正常，同时因为两项传感器略有分歧而中止起飞。这里“知道正常”与“应当继续”并不矛盾，因为行动依赖更丰富的命题集合。纯粹行动解释因而要求分析完整决策，而非把高风险本身当成魔法按钮。 反过来，风险有时改变证据取得方式。低风险时凭记忆填写地址，高风险合同时主动查阅产权文件；主体在复核前后拥有的证据实际上不同。若案例叙述说“证据完全相同”，却又让高风险主体注意到此前忽略的可能性，这个设定已经暗中改变了高阶证据。严谨分析要区分纯粹金钱后果、注意改变、获取新信息和既有证据的重新评价。

\minitopic{语言证据的边界}语境主义提出关于“知道”之语义的主张，不能只靠一个哲学家构造的案例决定。至少可采用四类语言证据：跨语境的撤回行为、嵌入句中的稳定性、说话者对过去归属的评价，以及不同表达之间的对照。若周五低风险时说“林知道”，周六高风险时又说“她从来不知道”，这支持标准回溯变化；若只说“现在不能靠这个办事”，则更支持行动解释。 量表问卷容易把语义、礼貌和道德评价混在一起。更可靠的实验会随机改变风险、证据和怀疑提示，要求参与者分别判断真值、断言得体性、责备程度与行动合理性，并测试效应是否由语言熟练度、专业背景或题目顺序造成。即使结果显示风险显著影响“知道”判断，也仍需排除参与者把“知道”理解成“确定无疑”或“足以行动”的语用强化。实验哲学能够约束理论，却不会自动替代语义论证。

\minitopic{比较视野：称谓、角色与时中}先秦材料通常不以知识归属句的语义为中心，而更关心言说者是否得其名、当其位、合其时。把这种角色敏感性直接翻译成当代语境主义并不准确：前者往往是关于实践分寸与礼的规范，后者首先是关于句子真值条件的理论。但两者可以形成有控制的比较问题：知识评价是否总能脱离说话场合和责任身份？一个判断在私人探索、课堂教学和公共决策中，是否承担不同的说明义务？ 比较研究应把“相似功能”与“相同概念”分开。儒家关于学、问、思、行的次序可提示认识活动嵌在角色关系中；庄子对立场与成心的反思可提示背景如何限定显著差异。\parencite{analects2003,zhuangzi2009} 但这些资源并不直接裁决“知道”的语义是否随归属者语境变化。其贡献在于扩大待解释现象：除了句子真值，还要解释谁有资格说、在何时说、说后承担何种行动后果。

\minitopic{综合案例：急诊分诊系统}某医院模型根据生命体征把患者标为“低风险”。在常规门诊，医生根据模型和检查结果让患者回家观察；在传染病暴发期间，同样分数触发隔离复检。请先画出信息结构：模型输出、医生看到的症状、模型在本院人群的校准数据、暴发这一背景信息、误放患者与误隔离患者的代价。然后分别提出三种解释。 语境主义版本关注不同谈话中“医生知道患者低风险”的归属标准；实践侵入版本主张暴发期的高代价提高知识所需证据；行动解释则说医生对个体风险的知识未必改变，但隔离决定受公共卫生理由约束。若暴发信息也提高患者属于高风险群体的基础概率，证据本身已经变化，不能再称为纯粹利害效应。若医生知道模型未在新毒株上验证，这构成高阶击败者。完整分析最后还要问：谁是归属者，谁承担决定，哪些复核可得，错误是否可逆？

\begin{tcolorbox}[breakable,colback=white,colframe=blue!30!black,title={专题研读任务},fonttitle=\bfseries]
将一个高风险案例改写成八个最小版本：分别改变归属者语境、主体利害、怀疑提示与证据质量。对每版独立判断知识、断言、行动和责备，再说明哪一项结果真正区分语境主义、实践侵入和纯粹行动解释。
\end{tcolorbox}

\index{语境主义}
\index{实践侵入}
\index{断言规范}
\index{知识归属}
\index{高风险}
